\section{形式的边界：生成物何时能被确认}

上一节留下的这三个判断，都指向同一个尚未展开的对象：智能输出来的东西究竟是什么。使用大模型时，人们最先遇到的疑问往往并不抽象。它为什么有时回答得像专家，有时又会编造事实？为什么能写出像样的证明，却在关键一步跳过去？为什么代码看起来合乎规范，执行时却调用了不存在的接口？为什么因果解释叙述顺滑，落到实验、工程或组织行动上却不成立？如果接入搜索、数据库、代码执行器、证明器和外部工具，幻觉是不是就会消失？

这些问题看似分散，其实指向同一个更深的问题：大模型生成的是形式，还是已经被确认的形式？论文标题、证明步骤、代码结构、因果解释、实验假设和行动方案，都可以以语言形式出现，模型能够生成它们，说明它从人类留下的痕迹中恢复了某些稳定结构；但结构被恢复出来，并不等于它已经指向真实对象、服从严格规则、通过执行检验，或能够进入现实并承担后果。大模型真正改变的，不只是答案从哪里来，而是形式可以在缺少同步确认的条件下被大量生成。形式生成与形式确认的分离，是理解大模型能力、幻觉和现实边界的入口。

\subsection{形式的双重身份与自我担保的不可能}

\subhead{形式的双重身份}

这也是为什么大模型的错误常常令人困惑。它的错误不是粗糙的，不是简单的乱码，不是完全没有结构，相反，它经常以非常完整的形式出现：不存在的论文有合理标题、作者、期刊和摘要；错误的证明有定理、引理、推导语气和符号秩序；不可运行的代码有良好风格、清晰函数名和完整注释；未经验证的医学建议有病理解释、剂量建议和风险说明。它不像一个无知者乱说，更像一个熟练掌握形式的人在没有确认的地方继续生成，问题因此不在于它是否拥有形式，而在于形式是否已经被对象、规则、执行和现实反馈确认。

所谓形式，并不只是符号、公理或公式，更底层地说，形式是能够降低不确定性的可迁移约束。一个句子如何继续，一个证明如何推进，一个程序如何调用，一个行动如何展开，一个失败如何解释，都不是任意发生的，只要某种结构能跨内容、跨语境、跨任务反复约束后续可能性，它就成为智能必须保留的形式。形式让经验不再停留在个案中，使判断可以迁移，使推理可以检查，也使行动可以预先组织。形式逻辑的意义正在这里：它没有发明推理，而是把推理的有效性从具体内容中抽离出来。三段论、命题逻辑、谓词逻辑和符号演算，都在说明同一件事，推理是否有效，不取决于对象是否相同，而取决于结构是否保持，只要结构保持，内容可以替换。形式由此成为智能最强的压缩方式之一，它用极低冗余保存高迁移能力，人类能够在不同对象之间迁移判断，能够把一类问题转化为另一类问题，能够把经验沉淀为规则，正是因为形式提供了这种可迁移的骨架。

这条路线在历史上走得很远。欧几里得的《几何原本》把散落的测量经验整理成公理与推演，使结论不再依赖某一次测量是否准确，而依赖链条是否成立。布尔在 1854 年把逻辑运算写成代数，弗雷格在 1879 年构造出概念文字，把“所有”“存在”“如果”这类结构纳入可运算的符号系统；到罗素与怀特海的三卷《数学原理》，公理化已经覆盖到算术本身。支撑这条路线的是一个很清楚的信念：只要公理足够明确、规则足够严格，证明就可以被机械检查，推理的有效性便不再取决于谁在推。压缩率也在这里达到极致，一条规则可以覆盖无穷多个实例，一个结构可以替换掉全部具体内容，而被压掉的内容并不需要被记住。

这条路线在二十世纪初被推到了纲领的高度。希尔伯特在 1900 年巴黎国际数学家大会上提出的二十三个问题中，第二个问的是算术系统的一致性能否证明；1928 年他在波洛尼亚再次回到这个问题，并把完备性与可判定性一并列为目标，其中可判定性就是后来所说的判定问题。纲领的设想很明确：把数学整理成一组清楚的公理和规则，使每一个数学命题原则上都可以被机械地判定。这个设想把形式的可迁移性推到了极限，也把形式的担保问题逼到了台面上。如果形式只能担保被它覆盖的部分，那么它能不能担保自己。

\subhead{形式不能自我担保}

形式化让推演可以独立检查，也给出了精确的通用边界。足够强、一致、可有效公理化的算术理论存在不可判定句，在第二不完备性定理的相应条件下不能证明自身形式化的一致性。\cite{godel1931} 停机问题排除了判定所有任意程序的一种通用算法；算法信息论则研究描述长度与不可压缩性。\cite{turing1936,chaitin1987} 这些结论限制不同对象，并不取消具体程序和证明的可靠检验。

这三次回答指向同一个结果：形式只能担保被它覆盖的部分，覆盖之外的部分必须由别的东西担保。这一点并不削弱形式，它只是给形式划定了位置。形式依然是智能最强的压缩方式之一，但压缩率越高，被压掉的部分越需要在别处找回来。哥德尔、图灵和 Chaitin 处理的是形式系统内部的边界，而现实系统还有另一层边界：形式在某个尺度上成立，不等于它在另一个尺度上成立。

复杂系统把这个提醒带回现实。真实世界不是一张已经写完的命题表，也不是给定输入便得到预期输出的简单机器。同一个系统，在不同尺度上呈现不同有效变量：微观尺度谈粒子、相互作用和涨落，宏观尺度谈温度、压力、流量和相变，生物尺度谈功能、稳态、适应和修复，社会尺度谈制度、策略、信任和反馈。尺度不是世界的附属标签，而是压缩世界的方式，一个形式在某个尺度上有效，并不意味着它可以无损迁移到另一个尺度。智能的错误常常不是完全没有形式，而是形式越界：把语言尺度上的因果叙述当成现实尺度上的可干预因果，把局部规律当成系统控制，把短期反馈当成长期目标，把单点准确率当成部署可靠性，把规则内推演当成规则外行动，都会产生误判。形式一旦脱离任务、介质、尺度、反馈和边界条件，就会从智能的骨架变成错误的来源。

形式越界在工程中很常见。仿真模型在训练工况内准确，换一种工况就失效；风控规则在历史数据上表现良好，遇到没有出现过的行为模式就误判；调度算法在标准货型上稳定，遇到非标货物就要人工介入。每一次失效都不是规则写错了，而是规则的适用范围被当成了世界的大小。判断一个形式在哪里有效，需要知道它是在什么数据、什么边界条件和什么失败代价下被建立起来的，而这些信息通常不在形式本身里面。

形式确认在工程中确实被做到了，但代价可以量出来。1994 年英特尔奔腾处理器的除法指令出现查表缺陷，公司为此计提了约四亿七千五百万美元的更换费用，这一事件推动了形式方法在硬件验证中的地位。软件一侧的极端案例是 seL4 微内核：2009 年公布的机器检查证明表明，约七千五百行 C 代码的实现与它的抽象规范一致，证明脚本约二十万行 Isabelle，投入约二十人年，并且只覆盖到 C 代码层，汇编与启动代码不在其中\cite{klein2009sel4}。编译器一侧的对应工作是 CompCert，一个用辅助证明工具验证过的 C 编译器，保证源代码的语义在编译之后仍然保持\cite{leroy2009compcert}。这些数字给出了形式确认的成本量级：被验证的代码规模越小、规范越窄，证明才越可能完成。一旦对象扩展到开放环境、硬件抖动、并发交互和人类使用，规范本身就先成了难以确定的部分。

\subsection{形式的压缩来源与低损失的假象}

\subhead{形式的压缩来源}

人工智能的发展也可以沿着这条线理解。早期符号 AI 抓住了智能中可表达、可组合、可检查的一面，规则、谓词、专家系统、定理证明、知识图谱和符号规划，都证明形式对于智能不可或缺，它的限制在于形式来源过于依赖人工书写，难以覆盖开放世界中的常识、语境、例外、身体、目标和后果。现实行动不是静态公理库的自然展开，符号系统擅长在明确边界内推演，却难以吸收人类世界中那些模糊、含混、隐含、身体化、历史化的经验结构。大模型的转折，在于它不再主要依靠人手写形式，而是从文本、代码、论文、法律、对话、证明、程序和实验记录中学习形式。语言不是词语的散堆，而是人类经验、推理、制度、失败、修正和行动记录的沉积，逻辑、分类、蕴含、否定、递归、因果、反事实、证明模板、程序结构、目标达成和失败修正，都反复存在于这些沉积中，模型为了降低预测误差，必须吸收这些低冗余、高迁移、强约束的关系骨架。

语料中的变量绑定、指代、函数调用、证明模式和目标关系具有可复用结构，学到它们有助于降低预测误差。训练是否学到、能够多稳定地调用，还须通过任务检验。能量、数据分布、目标与优化分别承担作用，耗散本身不保证网络获得逻辑或因果能力。信息处理的物理账本中，有一条明确下界：
\begin{equation}
E_{\min}=k_B T\ln 2
\end{equation}
这里 $k_B$ 是玻尔兹曼常数，取值 $1.38\times10^{-23}\,\mathrm{J/K}$；$T$ 是环境温度，单位为开尔文；$E_{\min}$ 是擦除一比特信息必须向环境排放的最小热量，单位为焦耳。在室温 $T\approx300\,\mathrm{K}$ 下，$E_{\min}\approx2.9\times10^{-21}\,\mathrm{J}$。这个下限来自熵的记账：擦除一比特把状态空间缩小一半，熵至少减少 $k_B\ln 2$，这部分熵必须由环境承接。该式给出的是下限而非实际值，现有计算设备的每次操作耗散远高于它，差距来自器件、存储层次、供电和散热结构，而不来自这个下限本身。它在这里的作用是提醒：信息处理不是无实体的符号操作，任何一次不可逆的写入或擦除都要在物理账本上支付代价。

但耗散只是必要条件，不是充分条件，并非所有耗散系统都会产生可迁移的逻辑结构。真正起作用的是几类条件叠加：电力、芯片和冷却系统提供真实能量通量；人类记录提供高度结构化的数据分布；损失函数把预测偏差转化为可比较的差异；优化算法把这种差异持续压入参数空间。耗散让计算得以发生，数据分布提供可压缩结构，损失函数提供选择压力，优化过程提供塑形机制，形式不是从能量中自动冒出，而是在能量驱动的优化过程中，被结构化数据和误差压力共同筛选出来。

这个图像不应被当作严格物理定律，但它有助于理解大模型为什么会长出语言、逻辑和代码结构：原始网络像尚未成形的地貌，训练数据像长期降下的暴雨，优化过程像持续冲刷，地貌材料当然有差别，但在足够大的数据、足够长的训练和足够强的误差压力下，真正决定形状的，是水流反复选择的通道。语言、逻辑、代码和因果叙述，正是穿过人类记录时最能降低阻力、最能排出误差、最稳定复现的路径，模型看起来懂语言和逻辑，不是因为它被灌入显式理性，而是因为这些结构是降低预测误差的高效通道。因此，现代大模型不是 if-else 规则集合，也不是传统符号规则库，它更像由物理耗散支撑、由信息梯度塑形、由参数空间承载的低冗余高约束结构，若用形象说法，它像一种被电力、数据和误差梯度共同“长”出来的信息晶体。但这里的“晶体”不是静止物体，而是需要持续能量输入才能被训练、调用和更新的动态结构，它的有序性不是免费的，而是由巨大外部耗散换来的；它的稳定性也不是绝对的，而是在训练分布、参数结构、推理上下文和确认机制之间被持续维持。

\subhead{低损失不等于真实}

边界也在这里出现。优化过程选择的是低损失通道，不是真理本身；预测训练逼出的是形式结构，不是已经被确认的形式。一个结构只要在语料中反复出现、能够降低预测误差，就可能被模型吸收，但它是否指向真实对象，是否服从严格规则，是否能被执行，是否能进入现实行动，还需要额外确认。\key{低损失不等于真实，流畅不等于有效，相容不等于可执行。}GPT-4 的技术报告就是一个很好的现实锚点：它展示了大模型在许多标准化考试和任务上的强大表现，也明确承认模型仍可能产生幻觉、推理错误，并在高风险场景中需要谨慎使用\cite{openai2023gpt4}。这个并置很重要，能力和边界不是互相否定的，同一个系统可以在考试、写作、推理、代码和多模态理解上表现惊人，同时仍然不能把自己的输出自动转化为已经确认的事实、已经验证的证明、已经可运行的程序或已经可部署的行动方案：它强大，正是因为它恢复了大量形式；它危险，也正是因为形式可以在未经确认时显得非常可信。

幻觉由此产生。它不是大模型能力之外的偶然故障，而是高维关联能力的阴影。高维关联能够把相似领域、相似格式、相似论证和相似目标连接起来，形成低阻力的生成路径，但如果这条路径没有被对象、规则和现实反馈截断，它就会继续延伸，生成结构相容却未经确认的内容。一个不存在的论文，可以由真实领域、真实术语、真实作者风格、真实期刊格式和真实论证套路组合出来；一条错误因果链，可以由相似场景、常见解释、局部相关性和叙述顺序连接出来；一个错误证明，可以保留证明的语气、步骤和符号外形，却在关键转折处失去规则约束。幻觉不是没有形式，而是形式生成超过了对象确认、规则检验和现实反馈。这也是为什么“希望模型完全没有幻觉”虽然自然，却不能只靠模型本体来实现：幻觉的根源不是单个事实没记牢，而是高维关联在缺少确认时仍会继续延伸，这个机制不能简单切掉，因为切掉它，也就切掉了模型的类比、补全、迁移和创造能力；真正可行的方向，不是让模型停止关联，而是让关联获得确认，哪些关联对应对象，哪些关联服从规则，哪些关联能经受执行，哪些关联能进入现实反馈。幻觉研究已经成为大模型可靠性研究中的核心问题，相关综述通常把幻觉理解为流畅、连贯但事实上错误或缺乏证据支持的生成内容，并把原因追溯到数据、训练目标、模型结构、推理过程和评估方式等多个环节\cite{huang2023hallucination}。这个现象提醒我们，幻觉不是“模型还不够大”这样一个简单问题，也不是“接上搜索就解决”这样一个局部工程问题，它是形式生成机制与确认机制不匹配的表现：大模型能够在形式空间中快速生成候选，但候选是否可靠，需要进入另一个过程。

\subsection{三种确认位置与确认的成本}

\subhead{三种确认位置}

证明是第二个重要案例。一个数学证明不是证明语气，不是“显然”“由此可得”“不难证明”这些连接词，也不是定理、引理、推论的外观，证明的核心是每一步变换都被规则允许，并且整个链条能被检查。AlphaGeometry 的价值正在这里：它在 30 道近年国际数学奥林匹克几何题中解出 25 道，接近人类金牌选手水平\cite{trinh2024alphageometry}，更重要的是，它不是单靠语言模型生成一段“像证明”的文本，而是把神经生成与符号演绎结合起来，神经部分提出有用构造，符号引擎负责推演和检查。这个例子说明，形式生成非常重要，但可靠证明必须接入规则确认，模型可以提出方向，证明必须被规则系统压实。代码是第三个案例。大模型写代码的能力已经足以改变软件生产方式，但代码也最清楚地显示了形式确认的必要性：代码看起来像代码，和代码能运行、能通过测试、能解决真实 issue、能在依赖环境中稳定工作，是不同层级的问题。SWE-bench Verified 包含 500 个经过人工筛选的真实 GitHub issue，用来评估模型或智能体能否在真实仓库中解决软件问题\cite{swebench2024}，这个基准的价值，不在于又多了一个排行榜，而在于它把代码从“文本生成”带回“执行确认”：仓库、依赖、测试、issue 描述、补丁、环境和回归验证共同构成确认链条，没有这条链条，代码只是形式完整的候选。

这三个案例指向同一个结构。论文引用需要对象确认，证明需要规则确认，代码需要执行确认，大模型可以在三者的外形上表现出惊人熟练，但外形不是确认：一个引用是否真实，需要检索和溯源；一个证明是否成立，需要逐步检查；一段代码是否可用，需要运行、测试和环境适配。形式生成可以在模型内部完成，形式确认却必须穿过外部对象、规则系统和现实执行。形式确认至少有三个位置。第一，表征确认：这个形式是否指向稳定对象？论文、作者、概念、接口、事实和关系，不能只看起来像真，还要能被检索、溯源和交叉核验。第二，推演确认：这个形式是否能按规则稳定变换？证明要能被检查，代码要能执行，计算要能复现，因果链要能承受反事实约束。第三，行动确认：这个形式是否能改变现实状态，并承受改变后的反馈？行动方案不只要语言完整，还要穿过接口、权限、资源、时间、失败恢复和责任边界。同一结构可以穿过这三个位置：因果可以作为文本图式被识别，可以作为推演结构被计算，也可以作为现实干预被验证，语言中成立，不等于推演中有效；推演中有效，不等于现实中可行；现实中可行，还要面对资源、时间、结构和责任。形式不是一次获得可靠性，而是在不同位置接受不同检验。Pearl 对因果的形式化工作正说明了这一点：因果不是把相关性讲得更顺，而是要引入干预和反事实结构，使某种依赖关系能够经受“如果改变条件会怎样”的检验\cite{pearl2009causality}。这三层落在第一章给出的验证金字塔的中段：那里从下到上排的是形式验证、语义验证、因果验证、工程验证、生命长期验证与社会责任验证，本章的表征确认、推演确认、行动确认分别对应其中的语义、因果与工程三层。大模型把候选生成成本降到极低，文本、代码、解释、假设和计划可以迅速涌出，但候选要变成可依赖结果，必须逐级上升：先确认对象，再确认规则，再确认行动，再往上，还要确认它能否嵌入系统，能否在长期运行中承受扰动、异常、审计和责任。

\subhead{确认的成本与边界}

训练阶段的耗散，把形式压缩出来；验证阶段的耗散，把形式变成可信。生成可以在参数空间中被预先摊薄，确认却必须在对象、规则、行动、系统和责任链中逐层支付。生成发生在金字塔底部，可靠性只能在上升过程中逐级获得。确认成本因此不是线性增加，在低风险、低耦合、可回滚的任务中，确认成本可能很低；在医疗、金融、法律、工程、公共政策和生命技术中，确认成本会迅速上升，越往上，确认越接近现实，也越接近资源、时间、组织、责任和承受。这个差别可以用通过率的形式写出来
\begin{equation}
N_{\text{cand}}=\frac{M_{\text{use}}}{\prod_{i=1}^{k}p_i}
\end{equation}
$k$为确认阶段数，$p_i$定义为通过前面各级之后，再通过第$i$级的条件概率，因此连乘不要求各级独立。$M_{\mathrm{use}}/\prod_i p_i$在稳定、可比较的重复尝试下给出期望或规划近似，不保证某一批一定成功。若使用无条件通过率，才需要额外独立性假设；相关性对估计的影响没有统一方向。模型改进也可能提高通过率、改善试验与减少返工。

检索、数据库、执行器、证明器、测试环境、API、传感器和反馈回路，都应放在这个位置理解：它们不是本节要综述的技术路线，而是不同形式的确认接口，检索确认对象，执行确认规则，测试确认边界，工具确认接口，反馈确认行动后果。它们能够降低幻觉，却不能消灭幻觉，因为确认路径本身也会出错：来源可能不可靠，测试可能覆盖不足，接口可能被误用，反馈可能延迟，责任可能无法分配。幻觉不会简单消失，只会从模型内部生成问题，转化为模型、工具与环境闭环中的确认问题。这也是为什么 RAG、工具调用、代码解释器和智能体框架不能被理解为“消除幻觉的按钮”：它们扩展了确认链条，但也带来了新的确认对象，检索结果是否真实，数据库是否过期，工具调用是否正确，执行环境是否覆盖真实条件，证明器处理的是不是已经形式化正确的问题，传感器反馈是否被正确解释；当模型开始调用外部世界，错误不再只是生成错误，而可能变成检索错误、接口错误、执行错误、覆盖错误、权限错误和责任错误，确认不是消失了，而是迁移了。由此可以得到一个可检验的边界声明：凡是声称模型已经获得通用推理能力或现实行动能力的判断，都必须接受三重检验，在没有外部确认接口的情况下，它能否稳定避免表征幻觉？能否持续遵守规则推演？能否完成可反馈、可回滚、可担责的现实行动？若不能，它拥有的就仍然主要是形式生成能力，而不是已经完成确认的智能能力。这个边界不会因为模型变大而自动消失，只会随着工具、环境和反馈系统的加入，被重新分配到更大的确认闭环之中。

这也解释了大模型与世界模型之间的距离。大模型从痕迹中学习形式，世界模型还要让形式接受现实反应的校准：语言中的因果可以由叙述维持，现实中的因果必须由干预检验；语言中的目标可以由句子表达，现实中的目标必须穿过资源、时间、结构和责任。预测有效的形式，还不是可干预世界的充分表达，它必须在行动后果中获得确认。一旦进入行动层，形式确认就和技术黏滞相接：确认不是免费的，事实确认需要检索、数据库和溯源，推演确认需要执行、测试和覆盖，行动确认需要传感器、接口、权限、试验、监测和回滚，系统确认需要流程改造、安全冗余、审计日志、责任分配和社会吸收，长期确认需要时间、异常、激励扭曲和代际后果。验证金字塔越往上，越会遇到资源、时间、结构、验证和责任的黏滞，大模型把候选生成推到极低成本，却没有取消候选变成可依赖结果所需支付的现实账本。同样是大模型应用，在不同介质中速度差异巨大：比特空间中，文本、代码、检索、总结和初步分析可以快速迭代，错误有时可以低成本撤回，候选可以低成本重写；原子空间中，行动需要材料、能量、传感器、接口、工况、测试和维护，错误可能造成损坏、停机或安全事故；生命空间中，确认更加昂贵，药物、基因编辑、发育过程、免疫反应、生态影响和代际后果都不能像文本一样反复撤销。形式越靠近现实深层，确认越重，黏滞越显影。

大模型的革命不是摆脱形式，而是改变了形式的生产机制。它让形式从人工规则进入高维表征，从显式书写进入统计学习，从少量公理进入大规模痕迹，但这场革命没有取消形式的边界。预测训练可以逼出结构，梯度冲刷可以塑造低损失通道，外部工具可以增加确认路径，却不能自动赋予形式以对象、规则、行动和责任上的可靠性；形式进入工具、环境和反馈之后，边界不是消失，而是被重新分配到系统闭环之中。所以，智能不只是生成形式，而是知道形式由什么确认、在什么尺度成立、何时需要被更换。大模型压缩的不是逻辑本身，而是逻辑、程序、因果和行动在人类痕迹中的形式，它的边界也不是没有形式，而是预测中形成的高维关联还没有自动获得表征、推演、行动、系统和长期检验下的可靠性。形式使智能可能，检验使形式进入现实。而一旦检验要发生，就先要有一个可以接受检验的地方，这要求模型与它的对象之间已经存在一条能够传递后果的通道。
